\documentclass[11pt,a4paper]{article}

% ------------------------------------------------------------------
% Packages
% ------------------------------------------------------------------
\usepackage[utf8]{inputenc}
\usepackage[T1]{fontenc}
\usepackage{lmodern}
\usepackage[margin=2.5cm]{geometry}
\usepackage{amsmath,amssymb,mathtools}
\usepackage{booktabs}
\usepackage{graphicx}
\usepackage{xcolor}
\usepackage{listings}
\usepackage{caption}
\usepackage[colorlinks=true,linkcolor=blue!60!black,citecolor=blue!60!black,urlcolor=blue!60!black]{hyperref}

% ------------------------------------------------------------------
% Macros
% ------------------------------------------------------------------
\newcommand{\E}{\mathbb{E}}
\newcommand{\Q}{\mathbb{Q}}
\newcommand{\dd}{\mathrm{d}}
\newcommand{\code}[1]{\texttt{#1}}
\DeclareMathOperator{\artanh}{artanh}

% Images extracted from Models/correlation.ipynb by make_figures.py ; a framed placeholder is printed while a file is missing.
\newcommand{\figfile}[2]{%
  \IfFileExists{figures/#1}{\includegraphics[width=#2]{figures/#1}}%
  {\fbox{\parbox[c][4cm][c]{#2}{\centering\small missing figure:\\ \texttt{figures/#1}}}}}

\lstset{basicstyle=\ttfamily\footnotesize, keywordstyle=\color{blue!70!black}, commentstyle=\color{gray},
        stringstyle=\color{teal}, frame=single, breaklines=true, language=Python, showstringspaces=false}

% ------------------------------------------------------------------
\title{Correlation Modelling for Worst-of and Best-of Options:\\
Statistical Tests on the LVMH--Hermès Dependence}
\author{Elias Benguedda}
\date{November 2026}

\begin{document}
\emergencystretch=3em
\maketitle

\begin{abstract}
Worst-of and best-of options pay on the minimum or the maximum of the performances of a basket, and their price
depends on the dependence between the underlyings as much as on their volatilities. This report prepares a
\emph{non-stochastic} correlation model for a basket of two French stocks, LVMH and Hermès: the pricing correlation
is a deterministic function of the history known at the pricing date, so that the same inputs always give the same
price. Before choosing that function, 21 years of daily data (5{,}588 returns, January 2005 to October 2026) are put
through a battery of tests. The correlation does not depend on the return horizon (0.60 from 1 to 60 days), so daily
returns can be used. It is not constant: equality across years is rejected ($p \approx 10^{-90}$), the yearly value
ranges from 0.10 to 0.85 with a level shift around 2015 (0.47 before, 0.74 after), and the Engle--Sheppard test
rejects a constant conditional correlation ($p \approx 10^{-5}$). On GARCH-standardised residuals, a mean-reverting
DCC model beats both the constant (likelihood ratio 443) and the integrated EWMA specification (likelihood ratio 60,
$p \approx 10^{-14}$); its persistence $a + b = 0.9915$ gives a half-life of 82 trading days, so that over a one-year
product only 41\% of today's deviation from the long-run correlation survives on average. The joint tails are fatter
than under a Gaussian dependence, mostly on the upside, and the correlation is higher in high-volatility regimes.
These results set the list of candidate models for the next sections. This paper is linked to my GitHub, which can be
found on my portfolio\footnote{\url{https://portfolio-eliasbenguedda.fr}}, and all scripts are available in the same
repository.
\end{abstract}

\tableofcontents
\newpage

% ==================================================================
\section{Introduction}\label{sec:intro}
% ==================================================================

\subsection{Motivation}
A worst-of (best-of) option on a basket of $n$ stocks pays on the worst (best) of the performances
$S^i_T / S^i_0$. These payoffs are the building blocks of most multi-asset structured products: worst-of
autocallables, worst-of reverse convertibles and best-of calls. Each vanilla option on a single underlying
identifies the volatility of that underlying, but no listed instrument identifies the correlation between two single
stocks with any liquidity. The pricing correlation therefore has to be estimated from history, and the estimator is
a modelling choice with a direct price impact.

A pricing desk also needs \emph{reproducibility}: a structured product priced twice with the same inputs must have
the same price, for instance when a data error forces a recalculation. This rules out correlation models whose
pricing correlation is itself random, such as a correlation following its own diffusion, and motivates a
\emph{non-stochastic} model:
\begin{equation}\label{eq:non-stochastic}
\rho_{\mathrm{pricing}}(t_0, T) = f\bigl((r_s)_{s \le t_0},\, T\bigr),
\end{equation}
a deterministic function of the returns observed up to the pricing date $t_0$, possibly depending on the maturity
$T$ of the product. The historical correlation over a window, its exponentially weighted version (EWMA) and the
forecast of a GARCH-type conditional correlation model are all of this form.

\subsection{Scope of this report}
Choosing $f$ in \eqref{eq:non-stochastic} requires knowing how the dependence between the underlyings actually
behaves. This first part of the study is purely statistical: on LVMH and Hermès, two luxury stocks listed in Paris,
it tests whether the correlation depends on the return horizon, whether it is constant, whether it mean-reverts and
at which speed, and whether a Gaussian dependence describes the joint extremes. Each test is chosen for its
consequence on the model; the models themselves and the pricing of worst-of and best-of options will be the subject of
the following sections.

\subsection{Organisation of the paper}
Section~\ref{sec:wobo} recalls the payoffs and how the correlation enters their price. Section~\ref{sec:data}
describes the data. Section~\ref{sec:marginals} tests the marginals, Section~\ref{sec:horizon} the dependence of the
correlation on the return horizon, Section~\ref{sec:constancy} its constancy, Section~\ref{sec:conditional} compares
constant, integrated and mean-reverting conditional correlation models, Section~\ref{sec:reversion} measures the mean
reversion without a model and Section~\ref{sec:tails} studies asymmetry and tails. Section~\ref{sec:synthesis}
draws the consequences for the model, Sections~\ref{sec:implementation} and~\ref{sec:limitations} describe the
implementation and the limitations, and Section~\ref{sec:conclusion} concludes.

% ==================================================================
\section{Worst-of and Best-of Options}\label{sec:wobo}
% ==================================================================

\subsection{Payoffs}
Let $X^i_T = S^i_T / S^i_0$ be the performance of stock $i$ and $m_T = \min_i X^i_T$, $M_T = \max_i X^i_T$. The
worst-of and best-of calls and puts of strike $K$ pay
\begin{equation}
(m_T - K)^+,\qquad (K - m_T)^+,\qquad (M_T - K)^+,\qquad (K - M_T)^+.
\end{equation}
For two assets, $m_T + M_T = X^1_T + X^2_T$ and the identity
$(m_T - K)^+ + (M_T - K)^+ = (X^1_T - K)^+ + (X^2_T - K)^+$ (and the same for puts) shows that the worst-of and the
best-of option of the same type add up to the two vanilla options, which do not depend on the correlation: whatever
the correlation adds to one, it takes from the other.

\subsection{Role of the correlation}
In a multi-asset Black--Scholes model, $\dd S^i_t / S^i_t = (r - q_i)\,\dd t + \sigma_i\,\dd W^i_t$ with
$\dd\langle W^i, W^j\rangle_t = \rho_{ij}\,\dd t$, the dependence between the terminal values is a Gaussian copula
with correlation matrix $(\rho_{ij})$. The spread between the two performances shrinks as the correlation rises:
$\E\,|X^1_T - X^2_T|$ decreases with $\rho$, so the worst performance rises and the best one falls. Worst-of calls are
therefore worth more, and worst-of puts less, when the correlation is high; best-of options move the other way. For
two assets, the call on the minimum and on the maximum have a closed form \cite{stulz1982}, which will serve to
validate the Monte Carlo pricer in the next part. A worst-of put, the option sold by the investor in a worst-of
autocallable, is \emph{short correlation} for the buyer of the note: a correlation estimate that is too high
underprices the risk.

The Gaussian copula has two consequences that the tests of this report check: the dependence is summarised by a
single number per pair, and the joint extremes are as rare as in a bivariate normal distribution (no tail dependence).

% ==================================================================
\section{Data}\label{sec:data}
% ==================================================================

\subsection{Source and adjustments}
Daily closes of LVMH (\code{MC.PA}) and Hermès (\code{RMS.PA}) on Euronext Paris are downloaded from Yahoo Finance
with the \code{yfinance} library \cite{yfinance} by the script \code{Data/Equity/download\_data.py}, from 3 January
2005 to 30 October 2026 (5{,}589 trading days, no missing value). Yahoo data is free but unofficial and its history
can be revised, so the files are frozen in the repository and every number of this report is reproducible.

Two series are stored per stock. The \emph{adjusted close}, corrected for dividends and splits, is used for every
statistic on returns, so that an ex-dividend date is not read as a price drop. The \emph{raw close} is the underlying
of the payoffs (price return); dividends will enter the pricing through the forwards. The log return is
$r^i_t = \ln(S^i_t / S^i_{t-1})$ on the adjusted close.

\subsection{Synchronicity}
Both stocks close at the same time on the same exchange, so their daily returns are synchronous. For pairs listed on
different time zones, a European close and an American close are separated by several hours of news, and the
correlation of daily returns is biased downwards; multi-day returns are then used. Section~\ref{sec:horizon} checks
that this is not needed here.

\begin{figure}[h]
\centering
\figfile{prices_returns.png}{\textwidth}
\caption{Total return of LVMH and Hermès since January 2005 (log scale) and daily log returns.}
\label{fig:prices}
\end{figure}

\begin{table}[h]
\centering
\begin{tabular}{@{}lccccccc@{}}
\toprule
 & Vol (ann.) & Skew. & Exc. kurt. & JB $p$ & LB $r$ $p$ & LB $r^2$ $p$ & ARCH-LM $p$ \\
\midrule
LVMH   & 28.3\% & 0.15 & 4.9 & $<10^{-300}$ & 0.018 & $10^{-154}$ & $6\times10^{-63}$ \\
Hermès & 27.3\% & 0.15 & 6.3 & $<10^{-300}$ & 0.23  & $4\times10^{-171}$ & $7\times10^{-89}$ \\
\bottomrule
\end{tabular}
\caption{Marginal statistics of the daily log returns, 2005--2026. JB: Jarque--Bera normality test; LB: Ljung--Box
test with 10 lags on returns and squared returns; ARCH-LM: Engle's test with 5 lags.}
\label{tab:marginals}
\end{table}

% ==================================================================
\section{Marginals}\label{sec:marginals}
% ==================================================================

\subsection{Tests}
Three properties of each return series are tested (Table~\ref{tab:marginals}): normality with the Jarque--Bera test
\cite{jarquebera1980}, absence of autocorrelation of the returns and of the squared returns with the Ljung--Box test
\cite{ljungbox1978}, and absence of conditional heteroskedasticity with Engle's ARCH-LM test \cite{engle1982}.

\subsection{Results}
Both stocks have an annualised volatility of about 28\%, a small positive skewness and an excess kurtosis of 5 to 6:
normality is rejected without ambiguity. The returns themselves are nearly unpredictable (Ljung--Box $p = 0.018$ for
LVMH, $0.23$ for Hermès), while the squared returns are strongly autocorrelated: volatility comes in clusters.

A GARCH(1,1) with Student-$t$ innovations \cite{bollerslev1986,bollerslev1987} is fitted to each stock
(Table~\ref{tab:garch}). The persistence $\alpha + \beta$ is 0.990 for LVMH and 0.996 for Hermès, and the degrees of
freedom of about 5 confirm fat-tailed innovations. The standardised residuals
$\varepsilon^i_t = (r^i_t - \mu_i) / \sigma_{i,t}$ remove the volatility clusters.

\begin{table}[h]
\centering
\begin{tabular}{@{}lccccc@{}}
\toprule
 & $\omega$ & $\alpha$ & $\beta$ & $\alpha + \beta$ & $\nu$ \\
\midrule
LVMH   & 0.037 & 0.050 & 0.939 & 0.990 & 5.1 \\
Hermès & 0.019 & 0.053 & 0.943 & 0.996 & 5.1 \\
\bottomrule
\end{tabular}
\caption{GARCH(1,1) with Student-$t$ innovations, returns in percent:
$\sigma^2_{t} = \omega + \alpha\,(r_{t-1} - \mu)^2 + \beta\,\sigma^2_{t-1}$.}
\label{tab:garch}
\end{table}

\subsection{Consequence for the dependence tests}
Volatility clustering matters for the correlation. When volatilities move through time, tests of the constancy of the
correlation applied to raw returns over-reject, and the sample correlation mechanically rises in volatile periods
\cite{forbesrigobon2002}. Every dependence test below is therefore run on the raw returns and, where it matters, on
the GARCH residuals $\varepsilon_t$. On the full sample the two correlations are close: 0.605 on raw returns and
0.609 on residuals.

% ==================================================================
\section{Correlation and Return Horizon}\label{sec:horizon}
% ==================================================================
The correlation of non-overlapping $h$-day log returns is computed for $h$ from 1 to 60 days, with the 95\% interval
given by the Fisher transform: $\artanh\hat\rho$ is approximately normal with variance $1/(n-3)$ \cite{fisher1921}.
A correlation that increases with the horizon would signal non-synchronous trading or a lead--lag effect (the Epps
effect, \cite{epps1979}).

\begin{table}[h]
\centering
\begin{tabular}{@{}rrccc@{}}
\toprule
Horizon (days) & $n$ & $\hat\rho$ & 95\% low & 95\% high \\
\midrule
1  & 5{,}588 & 0.605 & 0.588 & 0.621 \\
2  & 2{,}794 & 0.619 & 0.595 & 0.641 \\
5  & 1{,}117 & 0.599 & 0.560 & 0.636 \\
10 & 558     & 0.609 & 0.554 & 0.658 \\
20 & 279     & 0.556 & 0.470 & 0.632 \\
60 & 93      & 0.605 & 0.458 & 0.720 \\
\bottomrule
\end{tabular}
\caption{Correlation of non-overlapping $h$-day log returns.}
\label{tab:horizon}
\end{table}

The correlation is flat in the horizon, about 0.60 from 1 to 60 days, and every interval contains the daily value
(Table~\ref{tab:horizon}). There is no synchronicity bias: \emph{daily returns are used}, which gives five times
more observations than weekly returns and a correspondingly more precise estimator.

% ==================================================================
\section{Is the Correlation Constant?}\label{sec:constancy}
% ==================================================================

\subsection{Equality across years}
Split the sample into $K$ calendar years with $n_k$ observations and sample correlation $\hat\rho_k$. Under
independent normal observations, $z_k = \artanh\hat\rho_k$ is approximately $\mathcal N(\artanh\rho_k, 1/(n_k - 3))$
and, under $H_0: \rho_1 = \dots = \rho_K$,
\begin{equation}
Q = \sum_{k=1}^K (n_k - 3)(z_k - \bar z)^2 \sim \chi^2_{K-1}, \qquad
\bar z = \frac{\sum_k (n_k - 3)\, z_k}{\sum_k (n_k - 3)}.
\end{equation}
With $K = 22$ years, $Q = 490$ on raw returns ($p = 2\times10^{-90}$) and $Q = 429$ on GARCH residuals
($p = 10^{-77}$): removing the volatility clusters barely changes the conclusion. The yearly correlations
(Table~\ref{tab:yearly}) range from 0.10 in 2014 to 0.85 in 2022. The year 2014 is that of the settlement of LVMH's
stake in Hermès, which may explain an idiosyncratic decoupling; the low 2017 value (0.38) is not explained here.

\begin{table}[h]
\centering
\small
\begin{tabular}{@{}lccccccccccc@{}}
\toprule
Year & 2005 & 2006 & 2007 & 2008 & 2009 & 2010 & 2011 & 2012 & 2013 & 2014 & 2015 \\
\midrule
$\hat\rho$ returns   & 0.41 & 0.46 & 0.57 & 0.52 & 0.55 & 0.42 & 0.51 & 0.43 & 0.54 & 0.10 & 0.76 \\
$\hat\rho$ residuals & 0.41 & 0.39 & 0.55 & 0.57 & 0.53 & 0.49 & 0.50 & 0.44 & 0.52 & 0.23 & 0.74 \\
\midrule
Year & 2016 & 2017 & 2018 & 2019 & 2020 & 2021 & 2022 & 2023 & 2024 & 2025 & 2026 \\
\midrule
$\hat\rho$ returns   & 0.73 & 0.38 & 0.72 & 0.75 & 0.77 & 0.77 & 0.85 & 0.82 & 0.83 & 0.65 & 0.67 \\
$\hat\rho$ residuals & 0.70 & 0.39 & 0.69 & 0.73 & 0.78 & 0.77 & 0.85 & 0.81 & 0.84 & 0.67 & 0.68 \\
\bottomrule
\end{tabular}
\caption{Yearly correlations of daily returns and of GARCH residuals (2026: January to October).}
\label{tab:yearly}
\end{table}

\begin{figure}[h]
\centering
\figfile{rolling_correlation.png}{\textwidth}
\caption{Rolling correlation of daily returns over 63 and 252 days, yearly values and full-sample value.}
\label{fig:rolling}
\end{figure}

Figure~\ref{fig:rolling} shows a level shift around 2015: the correlation is 0.47 over 2005--2014 and 0.74 over
2015--2026, with fluctuations of $\pm 0.2$ around each level at the quarterly scale. The full-sample value of 0.61 is
representative of neither period. For a pricing correlation, the estimation window or the weighting of the
observations is therefore a first-order choice: a historical correlation over the whole sample would mix two
regimes. Over the last year the correlation is 0.67, over the last three years 0.70.

\subsection{Constancy of the conditional correlation}
The yearly test compares unconditional correlations. The test of Engle and Sheppard \cite{englesheppard2001}
addresses the conditional correlation directly. The standardised residuals are whitened with the constant
correlation matrix, $u_t = L^{-1}\varepsilon_t$ with $\bar R = LL^\top$ the Cholesky factorisation of the sample
correlation of $\varepsilon_t$. Under a constant conditional correlation, $u^1_t$ and $u^2_t$ are uncorrelated at
every date and the cross product $Y_t = u^1_t u^2_t$ has mean zero and no predictability. The regression
\begin{equation}
Y_t = \delta_0 + \delta_1 Y_{t-1} + \dots + \delta_s Y_{t-s} + \eta_t
\end{equation}
gives the statistic $\hat\delta^\top X^\top X \hat\delta / \hat\sigma^2_\eta \sim \chi^2_{s+1}$ under the null.

\begin{table}[h]
\centering
\begin{tabular}{@{}cccc@{}}
\toprule
Lags $s$ & Statistic & df & $p$-value \\
\midrule
1  & 23.3 & 2  & $9\times10^{-6}$ \\
5  & 33.1 & 6  & $1\times10^{-5}$ \\
10 & 47.1 & 11 & $2\times10^{-6}$ \\
\bottomrule
\end{tabular}
\caption{Engle--Sheppard test of constant conditional correlation on the GARCH residuals.}
\label{tab:es}
\end{table}

The constant conditional correlation is rejected at every lag (Table~\ref{tab:es}): past co-movements predict
future co-movements, which is what a dynamic correlation model exploits.

% ==================================================================
\section{Constant, Integrated or Mean-Reverting Correlation}\label{sec:conditional}
% ==================================================================

\subsection{Three specifications}
Following the two-step approach of Engle \cite{engle2002}, the conditional correlation is modelled on the GARCH
residuals through a pseudo-correlation matrix $Q_t$, with $\rho_t = Q_{12,t} / \sqrt{Q_{11,t}\,Q_{22,t}}$. Three
nested specifications are compared:
\begin{align}
\text{CCC \cite{bollerslev1990}:} \quad & Q_t = \bar Q, \\
\text{EWMA \cite{riskmetrics1996}:} \quad & Q_t = (1 - \lambda)\,\varepsilon_{t-1}\varepsilon_{t-1}^\top + \lambda\,Q_{t-1}, \\
\text{DCC \cite{engle2002}:} \quad & Q_t = (1 - a - b)\,\bar Q + a\,\varepsilon_{t-1}\varepsilon_{t-1}^\top + b\,Q_{t-1},
\end{align}
with $\bar Q$ the sample covariance of $\varepsilon_t$. The EWMA is the DCC on the boundary $a + b = 1$, without a
long-run target, and the CCC is the DCC with $a = 0$. All three are deterministic functions of the past and are
therefore candidates for \eqref{eq:non-stochastic}. The parameters maximise the correlation part of the Gaussian
log-likelihood,
\begin{equation}
\ell_C = -\frac12 \sum_{t} \left[\ln(1 - \rho_t^2) + \frac{(\varepsilon^1_t)^2 + (\varepsilon^2_t)^2
- 2\rho_t\,\varepsilon^1_t\varepsilon^2_t}{1 - \rho_t^2} - (\varepsilon^1_t)^2 - (\varepsilon^2_t)^2\right].
\end{equation}

\subsection{Forecasts over the life of a product}
The three models differ most in their forecasts. The CCC forecast is constant. The EWMA is integrated: its forecast
of $Q_{t+h}$ is $Q_{t+1}$ for every $h$, so the correlation of today is used for the whole life of the product. The
DCC reverts to the long-run correlation $\bar\rho$; with the approximation of Engle and Sheppard
\cite{englesheppard2001},
\begin{equation}\label{eq:dcc-forecast}
\E_t[\rho_{t+h}] \approx \bar\rho + (a + b)^{h-1}(\rho_{t+1} - \bar\rho),
\end{equation}
and the average correlation over a product of $H$ days is
\begin{equation}\label{eq:dcc-average}
\frac1H \sum_{h=1}^H \E_t[\rho_{t+h}] \approx \bar\rho + w(H)\,(\rho_{t+1} - \bar\rho), \qquad
w(H) = \frac{1 - (a+b)^H}{H\,(1 - a - b)}.
\end{equation}

\subsection{Results}
\begin{table}[h]
\centering
\begin{tabular}{@{}lccccl@{}}
\toprule
Model & $\ell_C$ & Parameters & AIC & BIC & Estimates \\
\midrule
CCC  & 1293.2 & 1 & $-2584$ & $-2578$ & $\bar\rho = 0.609$ \\
EWMA & 1484.8 & 1 & $-2968$ & $-2961$ & $\lambda = 0.977$ \\
DCC  & 1514.7 & 3 & $-3023$ & $-3004$ & $a = 0.031$, $b = 0.961$, $a + b = 0.9915$ \\
\bottomrule
\end{tabular}
\caption{Conditional correlation models fitted on the GARCH residuals (5{,}588 days). The DCC count includes the
targeted $\bar\rho$.}
\label{tab:dcc}
\end{table}

\begin{table}[h]
\centering
\begin{tabular}{@{}llccc@{}}
\toprule
$H_0$ & $H_1$ & LR & df & $p$-value \\
\midrule
CCC (constant)             & EWMA & 383.3 & 1 & $2\times10^{-85}$ \\
CCC (constant)             & DCC  & 443.1 & 2 & $6\times10^{-97}$ \\
EWMA (integrated, $a+b=1$) & DCC  & 59.8  & 1 & $1\times10^{-14}$ \\
\bottomrule
\end{tabular}
\caption{Likelihood ratio tests between the conditional correlation models.}
\label{tab:lr}
\end{table}

\begin{figure}[h]
\centering
\figfile{conditional_correlation.png}{\textwidth}
\caption{Conditional correlation of the GARCH residuals under the CCC, EWMA ($\lambda = 0.977$) and DCC models.}
\label{fig:conditional}
\end{figure}

Both dynamic models beat the constant correlation by a wide margin, and the DCC beats the EWMA
(Tables~\ref{tab:dcc} and~\ref{tab:lr}). The paths are close (Figure~\ref{fig:conditional}): the two models differ
less in their filtered correlation than in their forecast. The conditional correlation is \emph{persistent but
mean-reverting}: $a + b = 0.9915$ gives a half-life of $\ln 2 / |\ln(a+b)| = 82$ trading days, about four months.

Formula~\eqref{eq:dcc-average} turns this into weights on today's deviation from the long-run level:
\begin{center}
\begin{tabular}{@{}lccc@{}}
\toprule
Product maturity & 3 months & 1 year & 3 years \\
\midrule
Weight $w(H)$ of the current deviation & 0.78 & 0.41 & 0.16 \\
\bottomrule
\end{tabular}
\end{center}
Over a one-year product, only 41\% of the current deviation survives on average; over three years, 16\%. An EWMA,
whose forecast is flat, keeps 100\% at every maturity. The fitted EWMA parameter $\lambda = 0.977$ (an observation
half-life of 30 days) is also well above the 0.94 of RiskMetrics \cite{riskmetrics1996}, which was calibrated for
one-day risk forecasts, not for a correlation that must hold over the life of a product.

% ==================================================================
\section{Mean Reversion of the Realised Correlation}\label{sec:reversion}
% ==================================================================
The DCC result depends on the model. A model-free check uses the rolling correlation $\hat\rho^{(w)}_t$ over a
window of $w$ days. On its Fisher transform, the augmented Dickey--Fuller test \cite{dickeyfuller1979} tests the null
of a unit root (no mean reversion), and the first-order autocorrelation between consecutive non-overlapping windows
measures the persistence from one window to the next.

\begin{table}[h]
\centering
\begin{tabular}{@{}rcccccc@{}}
\toprule
Window (days) & Mean & Std & Min & Max & ADF $p$ & AR(1) non-overlapping \\
\midrule
63  & 0.61 & 0.21 & $-0.25$ & 0.92 & 0.0004 & 0.59 \\
126 & 0.61 & 0.19 & $-0.08$ & 0.90 & 0.054  & 0.60 \\
252 & 0.61 & 0.17 & $0.07$  & 0.89 & 0.13   & 0.70 \\
\bottomrule
\end{tabular}
\caption{Rolling correlation of daily returns: level, dispersion and mean reversion.}
\label{tab:reversion}
\end{table}

At the quarterly window, the unit root is rejected ($p = 0.0004$) and two consecutive quarters share an
autocorrelation of 0.59 only: short-term correlation reverts (Table~\ref{tab:reversion}). At the yearly window the
test loses power (overlapping windows, about 20 independent years) and the persistence is higher (0.70). The data
show two time scales: a slowly moving long-run level, which moved in 2015, and fast fluctuations around it. A pricing
correlation has to answer two questions: which long-run level, and how fast the current value reverts to it.

% ==================================================================
\section{Asymmetry and Tails}\label{sec:tails}
% ==================================================================
The multi-asset Black--Scholes model imposes a Gaussian copula. Three statistics test it, each compared to its value
under a bivariate normal distribution with the same correlation, since a Gaussian has non-trivial values of these
statistics too.

\subsection{Exceedance correlations}
Following Longin and Solnik \cite{longinsolnik2001}, the exceedance correlation at threshold $q$ is the correlation of
the standardised returns conditional on both being below $-q$ standard deviations (down) or above $q$ (up). For a
Gaussian it decreases with $q$ and is symmetric; the benchmark is computed by simulation ($2\times10^6$ draws).

\begin{table}[h]
\centering
\begin{tabular}{@{}cl rcc rcc@{}}
\toprule
 & & \multicolumn{3}{c}{Raw returns} & \multicolumn{3}{c}{GARCH residuals} \\
\cmidrule(lr){3-5}\cmidrule(lr){6-8}
$q$ (sd) & Side & $n$ & Data & Gaussian & $n$ & Data & Gaussian \\
\midrule
0   & down & 1{,}996 & 0.51 & 0.36 & 1{,}994 & 0.50 & 0.37 \\
    & up   & 2{,}029 & 0.51 & 0.36 & 2{,}029 & 0.44 & 0.36 \\
0.5 & down & 862 & 0.46 & 0.30 & 917 & 0.45 & 0.31 \\
    & up   & 819 & 0.52 & 0.30 & 898 & 0.42 & 0.31 \\
1   & down & 349 & 0.33 & 0.25 & 364 & 0.37 & 0.25 \\
    & up   & 288 & 0.45 & 0.25 & 325 & 0.42 & 0.26 \\
1.5 & down & 135 & 0.28 & 0.20 & 143 & 0.24 & 0.21 \\
    & up   & 119 & 0.36 & 0.22 & 109 & 0.47 & 0.22 \\
\bottomrule
\end{tabular}
\caption{Exceedance correlations against a Gaussian dependence with the same correlation.}
\label{tab:exceedance}
\end{table}

\begin{figure}[h]
\centering
\figfile{exceedance_correlation.png}{\textwidth}
\caption{Exceedance correlations (left of zero: both returns below $-q$; right: both above $q$).}
\label{fig:exceedance}
\end{figure}

The exceedance correlations are above the Gaussian benchmark at every threshold and on both sides
(Table~\ref{tab:exceedance}, Figure~\ref{fig:exceedance}). Part of the gap comes from the time variation of the
correlation itself: a mixture of Gaussians with different correlations has fatter joint tails than a single
Gaussian. On the GARCH residuals the up side stays high at 1.5 standard deviations (0.47 against 0.22) while the down
side comes back to the Gaussian level (0.24 against 0.21). For this pair there is no sign of a correlation that
rises in sell-offs; the strongest co-movements are joint rallies, consistent with sector-wide news for luxury stocks.

\subsection{Tail dependence}
On the ranks $U^i_t$ of the returns, the empirical lower tail dependence at level $q$ is
$\hat\lambda_L(q) = \#\{t : U^1_t < q,\, U^2_t < q\} / (nq)$, and $\hat\lambda_U(q)$ is defined in the same way on the
upper tail \cite{joe1997}. For a Gaussian copula $\lambda_L(q) \to 0$ as $q \to 0$.

\begin{table}[h]
\centering
\begin{tabular}{@{}c ccc ccc c@{}}
\toprule
 & \multicolumn{2}{c}{Raw returns} & & \multicolumn{2}{c}{GARCH residuals} & & \\
\cmidrule(lr){2-3}\cmidrule(lr){5-6}
$q$ & Lower & Upper & & Lower & Upper & & Gaussian \\
\midrule
10\% & 0.51 & 0.46 & & 0.50 & 0.41 & & 0.39 \\
5\%  & 0.41 & 0.39 & & 0.43 & 0.36 & & 0.31 \\
1\%  & 0.20 & 0.43 & & 0.30 & 0.27 & & 0.19 \\
\bottomrule
\end{tabular}
\caption{Empirical tail dependence against a Gaussian copula with correlation 0.605. At 1\% the lower tail of the raw
returns rests on 11 joint observations and the upper tail on 24.}
\label{tab:tail}
\end{table}

Both tails are above the Gaussian at 10\% and 5\% (Table~\ref{tab:tail}). At 1\% the estimates rest on a few dozen
joint observations at most and are not conclusive on their own; on raw returns the upper tail stands out (0.43
against 0.19), in line with the exceedance correlations.

\subsection{Volatility and market regimes}
The 21-day rolling correlation is averaged over the days in the top and bottom 20\% of the average individual
volatility, and over the days when the equally weighted basket fell or rose over the previous 21 days. Conditioning
on high volatility raises the sample correlation even when the true correlation is constant \cite{forbesrigobon2002};
the benchmark applies the same statistics to simulated returns with a constant correlation and the realised
volatility path of the data (average over 50 simulations).

\begin{table}[h]
\centering
\begin{tabular}{@{}lcc@{}}
\toprule
Regime & Data & Constant-correlation benchmark \\
\midrule
High volatility (top 20\%)    & 0.66 & 0.62 \\
Low volatility (bottom 20\%)  & 0.45 & 0.55 \\
Down market (21-day basket $<0$) & 0.64 & 0.59 \\
Up market (21-day basket $>0$)   & 0.58 & 0.59 \\
All days                         & 0.61 & 0.59 \\
\bottomrule
\end{tabular}
\caption{Mean 21-day rolling correlation by regime.}
\label{tab:regimes}
\end{table}

The correlation is 0.21 higher in high-volatility regimes than in low-volatility ones, against 0.08 for the benchmark
(Table~\ref{tab:regimes}): about 40\% of the gap is the conditioning bias, the rest is a genuine link between
correlation and volatility. Down markets show a slightly higher correlation than up markets (0.64 against 0.58, with
a benchmark of 0.59 on both sides). The correlation is thus linked to the volatility level, which matters for
worst-of products since their value is concentrated in the scenarios where the volatility is high.

% ==================================================================
\section{Synthesis and Consequences for the Model}\label{sec:synthesis}
% ==================================================================

\begin{table}[h]
\centering
\small
\begin{tabular}{@{}p{2.6cm}p{5.4cm}p{6.4cm}@{}}
\toprule
Test & Result & Consequence for the model \\
\midrule
Marginals & Fat tails, strong volatility clustering & Test the dependence on GARCH residuals; the volatility input matters as much as the correlation \\
Horizon & Correlation flat from 1 to 60 days & Daily returns, no synchronicity correction \\
Constancy & Rejected (Fisher, Engle--Sheppard); shift around 2015 & The window or the weighting is first-order; a 20-year average is biased \\
CCC / EWMA / DCC & DCC preferred; $a+b = 0.9915$, half-life 82 days; EWMA $\lambda = 0.977$ & A mean-reverting forecast over the product horizon is supported against a flat EWMA forecast \\
Mean reversion & Short-term correlation reverts; slow long-run level & Two scales: long-run level and reversion speed \\
Tails and regimes & Joint tails fatter than Gaussian, mostly upside; higher correlation in high volatility & A Gaussian copula understates joint extremes; correlation linked to the volatility level \\
\bottomrule
\end{tabular}
\caption{Summary of the tests.}
\label{tab:synthesis}
\end{table}

The tests (Table~\ref{tab:synthesis}) set the list of non-stochastic candidates for \eqref{eq:non-stochastic},
which the next part of the study will build and compare:
\begin{itemize}
\item \emph{Historical correlation} over a window (full sample, one year, three years, since the 2015 shift, or
equal to the maturity): the benchmark, whose dependence on the window measures the estimation risk;
\item \emph{EWMA}, on raw returns or on GARCH residuals, with $\lambda$ chosen by likelihood or by the forecast error
over the product horizon;
\item \emph{DCC forecast}: the average correlation over the life of the product, formula \eqref{eq:dcc-average},
which is the specification preferred by the likelihood and makes the pricing correlation depend on the maturity;
\item \emph{blend of a short-term and a long-term estimate} with maturity-dependent weights, a simple version of the
DCC forecast;
\item \emph{regime-switching forecast} \cite{hamilton1989}: the filtered regime probabilities give a deterministic
expected correlation, motivated by the 2015 shift and by the link with volatility;
\item \emph{Student-$t$ copula} instead of the Gaussian one: a fixed number of degrees of freedom adds tail
dependence, motivated by Section~\ref{sec:tails}.
\end{itemize}
The candidates will be ranked out of sample on two criteria: the error of the forecast against the correlation
realised over the life of the product, and the impact on the prices of worst-of and best-of options.

% ==================================================================
\section{Implementation}\label{sec:implementation}
% ==================================================================
\begin{itemize}
\item \code{Data/Equity/download\_data.py}: downloads the raw and adjusted closes and the dividends of LVMH and
Hermès from Yahoo Finance into \code{Data/Equity/};
\item \code{Models/correlation\_tests.py}: all the tests of this report (marginal tests, horizon, Fisher equality
test, Engle--Sheppard, CCC/EWMA/DCC fits and likelihood ratios, mean reversion, exceedance correlations, tail
dependence, regimes with their Gaussian or constant-correlation benchmarks). The recursions on $Q_t$ are linear
filters, so a likelihood evaluation over 5{,}588 days takes about a millisecond;
\item \code{Models/correlation.ipynb}: the notebook of this report;
\item \code{Reports/Correlation/make\_figures.py}: extracts the figures from the executed notebook.
\end{itemize}

\begin{lstlisting}
import pandas as pd
import correlation_tests as ct

raw = pd.read_csv("Data/Equity/LVMH_HERMES_Daily.csv", index_col="Date", parse_dates=True)
px = raw[["LVMH_AdjClose", "HERMES_AdjClose"]].set_axis(["LVMH", "HERMES"], axis=1)
r = ct.log_returns(px)
eps, garch = ct.garch_standardise(r)                    # GARCH(1,1)-t residuals
models, lr_tests, cond_rho, params = ct.compare_conditional_models(eps)
print(ct.engle_sheppard_test(eps), lr_tests, params["DCC"])
\end{lstlisting}

% ==================================================================
\section{Limitations}\label{sec:limitations}
% ==================================================================
\begin{itemize}
\item \emph{Two stocks, one period.} The conclusions hold for LVMH and Hermès over 2005--2026; the asymmetry of the
tails in particular (joint rallies rather than joint crashes) is specific to this pair and may not hold for a basket
of indices or of stocks from different sectors.
\item \emph{Assumptions of the tests.} The Fisher equality test assumes independent normal observations; it is
repeated on GARCH residuals, which removes the volatility clusters but not the fat tails. The likelihood ratio of CCC
against DCC has a parameter that is not identified under the null and the EWMA lies on the boundary of the DCC
parameter space, so the $\chi^2$ $p$-values are indicative; with ratios of 60 to 443 no correction changes the
conclusions.
\item \emph{Two-step estimation.} The correlation parameters are estimated given the GARCH fits, with a Gaussian
quasi-likelihood although the residuals are Student; the standard errors of the second step are not corrected for
the first.
\item \emph{In-sample.} All tests use the full sample. The model choice will be validated out of sample.
\item \emph{Data.} Yahoo Finance prices are free but unofficial; the adjusted close relies on Yahoo's dividend
adjustment.
\end{itemize}

% ==================================================================
\section{Conclusion}\label{sec:conclusion}
% ==================================================================
Over 21 years of daily data, the correlation between LVMH and Hermès averages 0.61 but is neither constant nor a
random walk. It shifted from about 0.47 to 0.74 around 2015, fluctuates by $\pm 0.2$ at the quarterly scale and
reverts with a half-life of about four months. A mean-reverting DCC model is preferred to both a constant correlation
and an integrated EWMA, and the joint extremes are more frequent than a Gaussian dependence implies, mostly on the
upside. For a non-stochastic pricing correlation, the main question is therefore less how to estimate the
correlation of today than how much of it to keep over the life of the product: with the fitted DCC, 41\% of today's
deviation over one year and 16\% over three years.

% ==================================================================
\appendix
\section{Notation}
% ==================================================================
\begin{center}
\begin{tabular}{@{}ll@{}}
\toprule
Symbol & Meaning \\
\midrule
$S^i_t$, $X^i_T = S^i_T / S^i_0$ & Price and performance of stock $i$ \\
$m_T$, $M_T$ & Worst and best performance \\
$r^i_t$ & Daily log return (adjusted close) \\
$\sigma_{i,t}$, $\varepsilon^i_t$ & GARCH conditional volatility and standardised residual \\
$\rho$, $\rho_t$, $\bar\rho$ & Correlation, conditional correlation, long-run correlation \\
$Q_t$, $\bar Q$ & Pseudo-correlation matrix of the DCC and its target \\
$\lambda$ & EWMA decay \\
$a$, $b$ & DCC innovation and persistence parameters \\
$w(H)$ & Weight of the current deviation in the average forecast over $H$ days \\
$\lambda_L(q)$, $\lambda_U(q)$ & Empirical lower and upper tail dependence at level $q$ \\
\bottomrule
\end{tabular}
\end{center}

% ==================================================================
\begin{thebibliography}{99}
\bibitem{yfinance} R. Aroussi, \emph{yfinance: download market data from Yahoo! Finance's API},
\url{https://github.com/ranaroussi/yfinance}.
\bibitem{stulz1982} R. M. Stulz, Options on the minimum or the maximum of two risky assets: analysis and
applications, \emph{Journal of Financial Economics}, 10(2):161--185, 1982.
\bibitem{jarquebera1980} C. M. Jarque and A. K. Bera, Efficient tests for normality, homoscedasticity and serial
independence of regression residuals, \emph{Economics Letters}, 6(3):255--259, 1980.
\bibitem{ljungbox1978} G. M. Ljung and G. E. P. Box, On a measure of lack of fit in time series models,
\emph{Biometrika}, 65(2):297--303, 1978.
\bibitem{engle1982} R. F. Engle, Autoregressive conditional heteroscedasticity with estimates of the variance of
United Kingdom inflation, \emph{Econometrica}, 50(4):987--1007, 1982.
\bibitem{bollerslev1986} T. Bollerslev, Generalized autoregressive conditional heteroskedasticity, \emph{Journal of
Econometrics}, 31(3):307--327, 1986.
\bibitem{bollerslev1987} T. Bollerslev, A conditionally heteroskedastic time series model for speculative prices and
rates of return, \emph{Review of Economics and Statistics}, 69(3):542--547, 1987.
\bibitem{bollerslev1990} T. Bollerslev, Modelling the coherence in short-run nominal exchange rates: a multivariate
generalized ARCH model, \emph{Review of Economics and Statistics}, 72(3):498--505, 1990.
\bibitem{fisher1921} R. A. Fisher, On the ``probable error'' of a coefficient of correlation deduced from a small
sample, \emph{Metron}, 1:3--32, 1921.
\bibitem{epps1979} T. W. Epps, Comovements in stock prices in the very short run, \emph{Journal of the American
Statistical Association}, 74(366):291--298, 1979.
\bibitem{englesheppard2001} R. F. Engle and K. Sheppard, Theoretical and empirical properties of dynamic conditional
correlation multivariate GARCH, NBER Working Paper 8554, 2001.
\bibitem{engle2002} R. F. Engle, Dynamic conditional correlation: a simple class of multivariate generalized
autoregressive conditional heteroskedasticity models, \emph{Journal of Business \& Economic Statistics},
20(3):339--350, 2002.
\bibitem{riskmetrics1996} J.P. Morgan and Reuters, \emph{RiskMetrics --- Technical Document}, 4th edition, 1996.
\bibitem{dickeyfuller1979} D. A. Dickey and W. A. Fuller, Distribution of the estimators for autoregressive time
series with a unit root, \emph{Journal of the American Statistical Association}, 74(366):427--431, 1979.
\bibitem{longinsolnik2001} F. Longin and B. Solnik, Extreme correlation of international equity markets,
\emph{Journal of Finance}, 56(2):649--676, 2001.
\bibitem{joe1997} H. Joe, \emph{Multivariate Models and Dependence Concepts}, Chapman and Hall, 1997.
\bibitem{forbesrigobon2002} K. J. Forbes and R. Rigobon, No contagion, only interdependence: measuring stock market
comovements, \emph{Journal of Finance}, 57(5):2223--2261, 2002.
\bibitem{hamilton1989} J. D. Hamilton, A new approach to the economic analysis of nonstationary time series and the
business cycle, \emph{Econometrica}, 57(2):357--384, 1989.
\end{thebibliography}

\end{document}
